\documentclass[11pt]{amsart}
\usepackage[margin=1in]{geometry}
\usepackage{amssymb,hyperref}
\newtheorem{theorem}{Theorem}
\newtheorem{proposition}[theorem]{Proposition}
\newtheorem{lemma}[theorem]{Lemma}
\theoremstyle{remark}
\newtheorem{remark}[theorem]{Remark}
\newcommand{\Z}{\mathbb{Z}}
\newcommand{\C}{\mathbb{C}}

\title[Norm-$m$ form of the alternating-product argument]{Remark R24: the alternating-product
argument of the circulant Hadamard proof for norm $m$, and three
parallel applications (quaternary, three-phase, Klein four-group)}
\author{[AUTHOR NAME]}
\date{9 October 2026}

\begin{document}
\begin{abstract}
OpenAI's preprint \emph{The circulant Hadamard conjecture} (family 179 of
\texttt{github.com/openai/math}) proves that an element $x\in\Z[i][C_{u^2}]$ with
$xx^*=u^2$ has an alternating product of character values that is a root of unity of
odd order, and combines this with a first-order $2$-adic comparison. The
proposition holds for every $x$ with $xx^*=u^2$, and squareness is not used in it.
We record that the same argument, with $u^2$ replaced by any $m$ and the comparison
done once, gives: no perfect quaternary sequence of length $2m$ ($m>1$ odd); no
perfect three-phase sequence of length $3m$ ($m>1$, $3\nmid m$); no perfect binary
array on $C_2\times C_2\times C_{u^2}$ ($u>1$ odd); and a rank-one identity for
$x\in\Z[i][C_{pq}]$ with $xx^*=pq$. All four are checked in Lean~4 on top of
the upstream formalization, using only the three standard axioms. \textbf{These are parallel
generalizations}: the method and almost all of the proof are the preprint's. We
claim no new idea.
\end{abstract}
\maketitle

\section{Statements}
Throughout, a function $a:G\to\C$ on a finite abelian group is \emph{perfect} if
$\sum_{g}a(g)\overline{a(g+s)}=0$ for all $s\neq0$. For $R\in\{\Z,\Z[i],\Z[\omega]\}$
and $N\ge1$ write $R[C_N]$ for the group ring, $x^*$ for the involution
$g\mapsto g^{-1}$ combined with complex conjugation on $R$, and $J=\sum_g g$.

\begin{proposition}[Prop$'$, Prop$''$]\label{prop}
Let $N>1$. Let $x\in\Z[i][C_N]$ with $xx^*=N$ and $N$ odd, or $x\in\Z[\omega][C_N]$
with $xx^*=N$ and $3\nmid N$. Then the alternating product
$\Delta(x)=\prod_{S}x_S^{(-1)^{|S|}}$ of~\cite[\S3]{OAI179} is a root of unity.
For each prime $p\mid N$ its order is a power of $p$, so the order is odd, resp.\ prime to $3$.
If $N$ has at least two distinct prime divisors, then $\Delta(x)=1$.
\end{proposition}

For $\Z[i]$ and $N=u^2$ this is \cite[Prop.~3.4, \texttt{prop:alternating}]{OAI179}.
The preprint itself says that the proposition ``applies to every $x\in\Z[i][P]$
satisfying $xx^*=u^2$'' (\texttt{introduction.tex} l.~191,
\texttt{contradiction.tex} l.~3). Its proof uses the norm only through the facts
that every character value has absolute value $\sqrt N$ and that $N$ is a unit away
from the primes dividing~$N$. Both facts hold for any $N$. Over $\Z[\omega]$, the localization layer
has to be redone at a prime above $3$ with $\lambda=1-\omega$, but the argument does not change.

\begin{theorem}[Master Lemma, sum form]\label{gml}
Let $(\mathcal O,\lambda,\ell)=(\Z[i],1+i,2)$ with $N>1$ odd, or
$(\Z[\omega],1-\omega,3)$ with $N>1$, $3\nmid N$. Let $x_0$ and $x_j$ ($j\in s$) lie in
$\mathcal O[C_N]$ with $xx^*=N$. Assume $x_j=x_0+t_jy_j$ with $0\neq t_j\in\lambda\mathcal O$, and
$\sum_j y_j=eJ+\lambda W$ with $e\equiv c\pmod\lambda$, $c\in\Z$, $\ell\nmid c$. Then this is
impossible.
\end{theorem}

For $\Z[i]$, $N=u^2$, $|s|=2$ this is the argument of \cite[\S4]{OAI179}.

\begin{theorem}\label{main}
\begin{enumerate}
\item[(Q)] For odd $m>1$ there is no perfect $a:\Z/2m\to\{\pm1,\pm i\}$.
\item[(T3)] For $m>1$, $3\nmid m$, there is no perfect $a:\Z/3m\to\{1,\omega,\omega^2\}$.
\item[(T1)] For odd $N>1$ there is no perfect $f:C_2\times C_2\times C_N\to\{\pm1\}$.
Equivalently, there is no Menon--Hadamard difference set in $C_2\times C_2\times C_{u^2}$, $u>1$ odd.
For non-square $N$ this is trivial, since $(\sum f)^2=4N$.
\item[(F)] Let $p\neq q$ be odd primes and $x\in\Z[i][C_{pq}]$ with $xx^*=pq$. Then
$\varepsilon(x)\,x(g)=\pi_p(x)(g\bmod p)\cdot\pi_q(x)(g\bmod q)$ for all $g$. Here $\varepsilon$ is
the augmentation and $\pi_p$ the push-forward to $C_p$. The statement is vacuous unless
$p\equiv q\equiv1\pmod 4$, since $\varepsilon(x)\overline{\varepsilon(x)}=pq$.
\end{enumerate}
\end{theorem}

\section{Why this is a parallel generalization}
The new ingredient for each alphabet is a one-line normalization. The rest is the
preprint's argument with $u^2$ replaced by $m$.
\begin{itemize}
\item (Q) Split $\Z/2m=C_2\times C_m$ and write $h=H_0+ZH_1$. Every element of
$\{\pm1,\pm i\}$ is $\equiv1\pmod{1+i}$, so $c=(H_0+H_1)/(1+i)$ and
$d=(H_0-H_1)/(1+i)$ lie in $\Z[i][C_m]$ and have norm $m$. Moreover $d=c+(i-1)H_1$ with
$H_1\equiv J\pmod{1+i}$. Theorem~\ref{gml} then applies with one index, $t=i-1$ and $e=1$. The
preprint needs two comparisons; here one suffices.
\item (T3) Write $h=H_0+ZH_1+Z^2H_2$ and $c_r=h(r)/(1-\omega)$ for $r\in\{1,\omega,\omega^2\}$.
Then $c_\omega-c_1=\lambda y_1$, $c_{\omega^2}-c_1=\lambda y_2$ and $y_1+y_2\equiv2\omega^2J$.
Apply Theorem~\ref{gml} over $\Z[\omega]$ with $e=2\omega^2$.
\item (T1) Write $h=\sum X^aY^bH_{ab}$ and $x_{\alpha\beta}=h(\alpha,\beta)/2\in\Z[C_N]$, which have norm $N$.
The three differences $x_{\alpha\beta}-x_{11}=-2b_j$ satisfy $\sum b_j\equiv J\pmod 2$.
Apply Theorem~\ref{gml} with three indices and $t=-2$.
\item (F) By Prop$'$, $\Delta=1$ for two primes. This reads
$\hat x(1)\hat x(\chi\psi)=\hat x(\chi)\hat x(\psi)$ for primitive $\chi$ on $C_p$ and $\psi$ on $C_q$.
Relabelling by units of $\Z/pq$ gives all character pairs, and Fourier inversion on
$C_p\times C_q$ finishes the proof.
\end{itemize}
By the gate used in our project (``could an expert write this down in an afternoon
from the original paper?''), the answer is yes. The 2-adic descent, the rigidity statements and the
classification of admissible normalizations are not on the proof path of Q, T1, T3
or F, and they are not part of this remark.

\section{Formal verification}
The Lean files in \texttt{scripts/lean/OAI/LinearAlgebra/CirculantHadamard/}
(\texttt{GeneralNorm/}: 7 files; \texttt{Eisenstein/}: 8 files; \texttt{RankOne/}: 2 files;
4085 lines) are added on top of \texttt{openai/math} commit
\texttt{fd4aeeb2ee4fc729c18d98444fed42fd0529eeeb}. They do not modify any upstream file
and import 15 unmodified upstream modules of the 179 formalization. The toolchain is Lean
\texttt{v4.34.1} with the upstream \texttt{lake-manifest.json} (Mathlib
\texttt{d13f23b723b8a846827a245b89c10fc7d3f11612}). \texttt{Checks.lean} prints the
statements and axioms of Prop$'$, Prop$''$, both Master Lemmas, Q, T1 (both forms),
T3 and F. Each depends only on \texttt{propext}, \texttt{Classical.choice} and
\texttt{Quot.sound}. The equivalence of perfect binary arrays with Menon--Hadamard
difference sets is standard and is not formalized.

\section{Computational sanity checks}
These are independent of the Lean development and share no code with the paper artifacts.
They are unit tests only; the theorems rest on the proofs.
Exhaustive search (normalized $a(0)=1$) finds no perfect quaternary sequence
of length $6,10,14$ and no perfect three-phase sequence of length $6,12,15$. As
positive controls it finds $2,8,32$ quaternary sequences of length $2,4,8$ and $6,54$
three-phase sequences of length $3,9$. A meet-in-the-middle over the four
normalized components finds no perfect binary array on $C_2\times C_2\times C_9$. On
$C_2\times C_2\times C_3\times C_3$ it finds $15552$, which is consistent with the known
$(36,15,6)$ Menon difference sets there~\cite{Menon}; one of them was checked directly.

\section{What is not claimed}
We claim no new method, and this remark makes no claim about individual lengths. Many lengths in
these families are excluded by self-conjugacy~\cite{Turyn,ADM,MaNg}, by field descent, or appear in
the tables of~\cite{Arasu}. For T1, Turyn's exponent bound already excludes
$C_2\times C_2\times C_{u^2}$ whenever some prime $p\mid u$, $p^a\parallel u$, is self-conjugate
modulo $2(u/p^a)^2$. This covers all prime powers $u$ and all odd $u<39$. Theorem~(F) is proved only for two
distinct odd primes and $\Z[i]$. Correctness depends on the Lean kernel and on the
faithfulness of the Lean statements; the reduction from difference sets to perfect arrays is not formalized.
Related work by the same authors is in preparation.

\section*{Acknowledgement and AI disclosure}
The method is OpenAI's~\cite{OAI179}. Developed with AI assistance (Anthropic Claude,
Opus 5.5) under [AUTHOR NAME]'s direction; independent adversarial reviews were run
by separate AI sessions and are not human peer review.

\begin{thebibliography}{9}
\bibitem{OAI179} OpenAI, \emph{The circulant Hadamard conjecture}, preprint,
September 23, 2026, \url{https://github.com/openai/math}, commit
\texttt{fd4aeeb2ee4fc729c18d98444fed42fd0529eeeb},
\texttt{preprints/The-circulant-Hadamard-conjecture-September-23-2026/}; Lean:
\texttt{lean/OAI/LinearAlgebra/CirculantHadamard/}.
\bibitem{Turyn} R.~J. Turyn, Character sums and difference sets, \emph{Pacific J. Math.} 15 (1965), 319--346.
\bibitem{Arasu} K.~T. Arasu, Sequences and arrays with desirable correlation properties, in
\emph{Information Security, Coding Theory and Related Combinatorics}, IOS Press, 2011, 136--171.
\bibitem{ADM} K.~T. Arasu, W.~de~Launey and S.~L. Ma, On circulant complex Hadamard matrices,
\emph{Des. Codes Cryptogr.} 25 (2002), 123--142.
\bibitem{MaNg} S.~L. Ma and W.~S. Ng, On non-existence of perfect and nearly perfect sequences,
\emph{Int. J. Inf. Coding Theory} 1 (2009), 15--38.
\bibitem{Menon} P.~K. Menon, On difference sets whose parameters satisfy a certain relation,
\emph{Proc. Amer. Math. Soc.} 13 (1962), 739--745.
\end{thebibliography}
\end{document}
